\section{Related Work}
\label{sec:related}

\noindent\textbf{BPE as an optimisation problem.}
BPE was introduced for data compression \citep{gage1994} and adapted to subword
segmentation by \citet{sennrich2016}. \citet{zouhar2023} formalise it as the search for a
merge sequence that maximises the number of tokens saved and bound the approximation ratio
of the greedy algorithm; Unigram tokenization \citep{kudo2018} prunes a large vocabulary
under a likelihood objective. Both objectives are defined on the pooled corpus, with no
notion of cost per language.

\noindent\textbf{Cross-lingual disparity and fair tokenizers.}
\citet{petrov2023} show that the tokenizers of deployed models encode the same content with
very different numbers of tokens across languages, measured on parallel text, the unit we
adopt. Parity-aware BPE \citep{foroutan2026} changes the merge rule: at every step it merges
the most frequent pair of the language with the worst compression on a multi-parallel
development set, and a hybrid schedule and a moving-window variant reach intermediate
tokenizers. \citet{kanjirangat2026} extend the parity rule to the $N$ worst-off languages,
and \citet{lee2026} find that Parity-aware BPE lies on the efficiency--equity frontier for
Southeast Asian languages but becomes less equitable as the vocabulary grows. Other work
changes the input rather than the rule, by rebalancing the data \citep{chung2023unimax} or
by a morphology-driven byte encoding \citep{limisiewicz2024myte}. In all of these, the
intermediate tokenizers come from schedules, a choice of $N$ or data weights, none derived
from an objective over per-language costs.

\noindent\textbf{Vocabulary allocation and Indian languages.}
VoCap \citep{zheng2021vocap} allocates vocabulary per language from its average log
probability, XLM-V \citep{liang2023xlmv} builds per-cluster vocabularies before merging
them, and \citet{brahma2026} set each script cluster's share for 17 Indian languages in
proportion to its share of the union of monolingual vocabularies. SuTRA \citep{rathore2026}
penalises merges that cross morpheme boundaries in Hindi, Marathi and Gujarati, whereas
\citet{arnett2025} find across 70 languages that morphological alignment explains little of
the variance in downstream performance. We use the SuTRA languages as our testbed.

\noindent\textbf{Fairness through scalarisation.}
Power means are a standard scalarisation in multi-objective optimisation
\citep{miettinen1999}; the same family appears as $\alpha$-fair utility in resource
allocation \citep{mo2000fair} and as the Atkinson index in economics \citep{atkinson1970}.
Our trainer builds on the lazy greedy technique of \citet{minoux1978}. To our knowledge,
this family has not been applied to the merge rule of a tokenizer, where it yields a
closed-form weighting of pair counts and an exact lazy trainer.
